\scanpage{89}
La théorie générale donne un iso (pas extérieur !)
\[
\tag{95}
\begin{aligned}
\Pi_1(X_n^*,\inserted{D_n;}P)
&\overset\sim\longrightarrow
\Pi_1(\underset{\overset{\mathrm{def}}=X_0\setminus\{0\}}{\underline{X_0^{\unsure{*|**}}}},2\{1\}+n\{\infty\};P_+)\\
&\simeq\Pi_{0,3}/\langle\rho_1^2,\rho_\infty^n\rangle,
\end{aligned}
\]
pour obtenir un isomorphisme précisé (93), il faut composer (95) avec un isomorphisme
\[
\tag{96}
\Pi_1(X_n^*,\inserted{D_n;}P_+)\overset\sim\longrightarrow
\Pi_1(X_n^*,D_n,P)
\]
déduit d'une classe de chemins $P_+\mathbin{-}P$ dans $X_n^*$ — on prendra la \inserted{l'unique} classe de chemins provenant d'un chemin $\alpha$ \uline{dans $r_0$}, comme choix le plus naturel, ce qui précise 96. L'homomorphisme (95), quant à lui, est donné par la formule des plus évidentes
\[
(l,g)\longmapsto g_n(l)
\]
où $g\in D_n$, $\underbrace{l:g(P_{\unsure{*|+}})\longrightarrow P_{\unsure{*|+}}}_{\substack{\text{chemin dans}\\X_n^*,\text{ ou mieux,}\\\text{chemin dans }X_n^{\unsure{***|**}}\\
=X_n^*\setminus\{(Q_a),P_\omega,P_{\omega'}\}\\
=X_n\setminus\{\text{pts de ramification}\}\\
=X_n\mid U_{03}}}$,
d'où $g_n(l):g_n(P)\longrightarrow g_n(P)$ un lacet en $P_+=g_n(P)$ dont on prend la classe dans le quotient de $\Pi_{03}$ — qui ne dépend pas du choix fait, mais seulement de l'élément de $\Pi_1(X_n^*,D_n;P)$ représenté par $(l,g)$ --.

Considérons donc l'application composée de (95)(96), qui \deleted{\unsure{$g_n$}} \inserted{non extérieure canonique}
\[
\tag{97}
\begin{aligned}
\Pi_1(X_n^*,D_n;P_+)&\xrightarrow{\deleted{\unsure{g}}\ g_n}
\Pi_1(X_0,\infty\{0\}+2.\{1\}+n\{\infty\};P_+)\\
&\simeq\Pi_{03}/\langle\rho_1^2,\rho_\infty^n\rangle
\end{aligned}
\]
et calculons ses valeurs pour $\rho,\sigma_{a_0}$. L'image de $\rho$ par 96, est donné par le chemin
\[
\rho(P)\xrightarrow{\alpha_\rho}\deleted{\unsure{P_+}}P^{\unsure{*|+}}
\quad\text{dans }X_n^{\unsure{***|**}}
\]
que fournit le 1-squelette de la subdivision

\margin{NB $\alpha_\rho$ est obtenu par déformation dans $X_n^{\unsure{*|**}}$ du chemin composé $\alpha\circ\rho(\alpha)^{-1}$.}

\margin{88}
